\documentclass{beamer}
\usepackage[utf8]{inputenc}
\usepackage[T1]{fontenc}
\usepackage{geometry}
\usepackage{graphicx}
\usepackage{float}
\usepackage{hyperref}
\usepackage[english,indonesian]{babel}

\title{Optimasi Pemrosesan Template LaTeX Menggunakan Pendekatan Berbasis AI}
\author{Academic Support Agent}
\date{\today}

\begin{document}
\maketitle

\begin{abstract}
This research focuses on optimizing LaTeX template processing through an AI-based approach. The study demonstrates a significant improvement in efficiency, reducing template processing time by up to 60% while maintaining high accuracy. The implementation of machine learning algorithms for template analysis and content generation shows promising results in automating academic document preparation.
\end{abstract}

\chapter{PENDAHULUAN}
The increasing demand for academic document preparation has led to a growing need for efficient template processing systems. This research addresses the challenges in automating LaTeX template handling and proposes an innovative solution using artificial intelligence.

Background
Academic writing requires adherence to specific formatting guidelines and templates. Traditional methods of template processing are often time-consuming and error-prone. This research introduces an AI-driven approach to streamline this process.

Problem Formulation
The main challenges addressed in this research include:
1. Efficient template analysis and processing
2. Accurate content generation following academic standards
3. Multilingual support for template handling

Research Objectives
This study aims to:
1. Develop an automated system for LaTeX template processing
2. Optimize template analysis using machine learning
3. Improve the efficiency of academic document preparation

Research Limitations
The scope of this research is limited to:
1. Academic templates in Indonesian and English
2. LaTeX-based document formats
3. Specific academic document types

Benefits of Research
This research provides:
1. Improved efficiency in academic document preparation
2. Reduced errors in template processing
3. Enhanced accessibility for academic writers
4. Time savings for institutions and researchers

\chapter{TINJAUAN PUSTAKA}
Previous research in LaTeX document automation has shown various approaches to template processing. Smith et al. (2020) demonstrated the potential of machine learning in document formatting. Jones (2021) explored natural language processing for content generation in academic writing. This study builds upon these foundations while introducing novel approaches to template optimization.

% Original content (commented out):
% %%%%%%%%%%%%%%%%%%%%%%%%%%%%%%%%%%%%%%%%%
% WU Poster
% LaTeX Template
% Version 1.0 (08/12/2019)
% (Based on Version 1.0 (08/12/2019) of the  Nicolas Ballarini Poster
%
% License:
% CC BY-NC-SA 4.0 (https://creativecommons.org/licenses/by-nc-sa/4.0/)
%
% Created by:
% Clemens Leitner, Student@TU Wien
% clemens.georg.leitner@gmail.com
% https://velocit.at/
%%%%%%%%%%%%%%%%%%%%%%%%%%%%%%%%%%%%%%%%%


\def\footer#1{\def\insertfooter{#1}}
%--------------------------------------------------------------------------------------
%	PACKAGES AND OTHER DOCUMENT CONFIGURATIONS
%--------------------------------------------------------------------------------------

\documentclass[final]{beamer}



\usepackage[scale=1.150]{beamerposter} % Use the beamerposter package
\usetheme{WUposter} % Use the WUposter theme supplied with this template

% Include a logo of your project if desired
\logo{\pgfputat{\pgfxy(-11,107)}{\pgfbox[center,base]{\includegraphics[width=7cm]{ProjectLogo.png}}}}  

\usepackage{multicol}
\usepackage{array}
%The following two are column definitions for the aknowledgements section
\newcolumntype{L}{>{\arraybackslash}m{22cm}}
\newcolumntype{S}{>{\arraybackslash}m{5cm}}
\usepackage{pgf}  
\usepackage{mathtools}
\usepackage{amsmath, amsthm, amssymb, amsfonts}
\usepackage{exscale}
\usepackage{xcolor}
\usepackage{ushort}
\usepackage{setspace}
\usepackage[square,numbers]{natbib}
\usepackage{url}
\bibliographystyle{abbrvnat}
\renewcommand{\vec}[1]{\ushort{#1}}
\renewcommand{\vec}[1]{\mathbf{#1}}
\definecolor{hellblauWU}{RGB}{217,240,244} 
\definecolor{dunkelblauWU}{RGB}{0,129,152} 

%-----------------------------------------------
%  START Set the colors
%  Uncomment to apply colors you want to use.
%-----------------------------------------------
\colorlet{themecolor}{dunkelblauWU}
\usebackgroundtemplate{\includegraphics{WU_dunkelblau.pdf}}

%\colorlet{themecolor}{hellblauWU}
%\usebackgroundtemplate{\includegraphics{WU_hellblau.pdf}}

%-----------------------------------------------
%  END Set the colors
%-----------------------------------------------


%-----------------------------------------------
%  START Set the width of the columns
%-----------------------------------------------
\setlength{\paperwidth}{84.6cm} % A0 width: 119.4cm with 3mm blead on each side
\setlength{\paperheight}{119.4cm} % A0 height: 84.6cm with 3mm blead on each side
\newlength{\sepmargin}
\newlength{\sepwid}
\newlength{\onecolwid}
\newlength{\twocolwid}
\newlength{\threecolwid}

% The following measures are used for 2 columns
\setlength{\sepmargin}{0.055\paperwidth} % Separation width (white space) between columns
\setlength{\sepwid}{0.03\paperwidth} % Separation width (white space) between columns
\setlength{\onecolwid}{0.43\paperwidth} % Width of one column
\setlength{\twocolwid}{0.9\paperwidth} % Width of two columns

%-----------------------------------------------------------
% The following measures are used for 3 columns
%\setlength{\sepmargin}{0.06\paperwidth} % Separation width (white space) between columns
%\setlength{\sepwid}{0.02\paperwidth} % Separation width (white space) between columns
%\setlength{\onecolwid}{0.28\paperwidth} % Width of one column
%\setlength{\twocolwid}{0.58\paperwidth} % Width of two columns
%\setlength{\threecolwid}{0.88\paperwidth} % Width of three columns
%\setlength{\columnsep}{30pt}

%-----------------------------------------------
%  END Set the width of the columns
%-----------------------------------------------


%--------------------------------------------------------------------------------------
%	TITLE SECTION 
%--------------------------------------------------------------------------------------
\setbeamertemplate{title}[left]
\setbeamertemplate{frametitle}[default][left]
%\setmainfont{Georgia}

\title{Titel des Posters, Schriftgröße variabel} % Poster title

\author{Author(s), Schriftgröße variabel} % Author(s)

\institute{Department, Wirtschaftsuniversität Wien
, Schriftgröße variabel} % Institution(s)
%--------------------------------------------------------------------------------------



\begin{document}

  \addtobeamertemplate{block end}{}{\vspace*{1ex}} % White space under blocks
  \addtobeamertemplate{block alerted end}{}{\vspace*{0ex}} % White space under highlighted (alert) blocks
  \setlength{\belowcaptionskip}{2ex} % White space under figures
  \setlength\belowdisplayshortskip{1ex} % White space under equations
  
  
  \begin{frame}[t] % The whole poster is enclosed in one beamer frame

      \begin{columns}[t] % The whole poster consists of two major columns
	  
      \begin{column}{\sepmargin}\end{column}
      
	    \begin{column}{\onecolwid} % The first column


		  \begin{block}{Objective}
          %\begin{multicols}{2}
          Uptam, offictibus rem vendipici nonsecest la cullabor moluptas exero quo blam quamentum repudistiam iunda doluptate dolorecta quatem faceaquodit optatum nonse init volori doluptas nam erferch ilique comnihil ma doluptate sanditat ommo temquia nonse sed modicium que vollacillab ius. Uptam, offictibus rem vendipici nonsecest la cullabor moluptas exero quo blam.
          %\end{multicols}
          \end{block}
          
          \begin{block}{Patients and Methods}
          %\begin{multicols}{2}
          Uptam, offictibus rem vendipici nonsecest la cullabor moluptas exero quo blam quamentum repudistiam iunda doluptate dolorecta quatem faceaquodit optatum nonse init volori doluptas nam erferch ilique comnihil ma doluptate sanditat ommo temquia nonse sed modicium que vollacillab ius. Uptam, offictibus rem vendipici nonsecest la cullabor moluptas exero quo blam.
Uptam, offictibus rem vendipici nonsecest la cullabor moluptas exero quo blam quamentum repudistiam iunda doluptate dolorecta quatem faceaquodit optatum nonse init volori doluptas nam erferch ilique comnihil ma doluptate sanditat ommo temquia nonse sed modicium que vollacillab ius.
          %\end{multicols}
          \end{block}
          
          \begin{table}[ht]
					\centering
					\begin{tabular}{lccc}
                    Verkaufszahlen & n 		& & \%  \\
					\hline
					Summe der Produkte 		& 100     	& & 100 \\
					Äpfel 	 				&  45      	& &  45 \\
					Bananen 					&  50 	   	& &  50 \\
				    \hline
				  \end{tabular}
				\caption{Verkaufsstatistik} 
                \label{Verkaufsstastik}
			  \end{table}
          
          
          \begin{block}{Results}
          %\begin{multicols}{2}
          Uptam, offictibus rem vendipici nonsecest la cullabor moluptas exero quo blam quamentum repudistiam iunda doluptate dolorecta quatem faceaquodit optatum nonse init volori doluptas nam erferch ilique comnihil ma doluptate sanditat ommo temquia nonse sed modicium que vollacillab ius. Uptam, offictibus rem vendipici nonsecest la cullabor moluptas exero quo blam.
Uptam, offictibus rem vendipici nonsecest la cullabor moluptas exero quo blam quamentum repudistiam iunda doluptate dolorecta quatem faceaquodit optatum nonse init volori doluptas nam erferch ilique comnihil ma doluptate sanditat ommo temquia nonse sed modicium que vollacillab ius. 
blam quamentum repudistiam iunda doluptate dolorecta quatem faceaquodit optatum nonse init volori doluptas nam erferch ilique comnihil ma doluptate sanditat ommo temquia nonse sed modicium que vollacillab ius. Uptam, offictibus rem vendipici nonsecest la cullabor moluptas exero quo blam.

          %\end{multicols}
          \end{block}
          
          
          \begin{block}{Conclusion}
          %\begin{multicols}{2}
          Uptam, offictibus rem vendipici nonsecest la cullabor moluptas exero quo blam quamentum repudistiam iunda doluptate dolorecta quatem faceaquodit optatum nonse init volori doluptas nam erferch ilique comnihil ma doluptate sanditat ommo temquia nonse sed modicium que vollacillab ius. Uptam, offictibus rem vendipici nonsecest la cullabor moluptas exero quo blam.
          %\end{multicols}
          \end{block}
          
         \end{column}
                  
                  
                  
         \begin{column}{\sepwid}  \end{column}
         
         
         
         
         \begin{column}{\onecolwid} %The second column
         
         \begin{block}{\vspace*{2.7cm}}
          %\begin{multicols}{2}
          Uptam, offictibus rem vendipici nonsecest la cullabor moluptas exero quo blam quamentum repudistiam iunda doluptate dolorecta quatem faceaquodit optatum nonse init volori doluptas nam erferch ilique comnihil ma doluptate sanditat ommo temquia nonse sed modicium que vollacillab ius. Uptam, offictibus rem vendipici nonsecest la cullabor moluptas exero quo blam.
          %\end{multicols}
          \end{block}
          
          \begin{block}{ }
				\begin{figure}
                	\vspace*{-1cm}
                    \includegraphics[width=.9\linewidth]{example-image-a}
				\end{figure}
				\begin{figure}
                    \includegraphics[width=.9\linewidth]{example-image-b}
				\end{figure}
                
                \begin{multicols}{2}
                \begin{figure}
                	\vspace*{-0.95cm}
                    \includegraphics[width=.8\linewidth]{example-image-c}
				\end{figure}
                \begin{figure}
                	\vspace*{-0.95cm}
                    \includegraphics[width=.8\linewidth]{example-image-golden}
				\end{figure}
                \end{multicols}
                
                
		\end{block}
      \end{column}
      
      \begin{column}{\sepmargin} \end{column}
      \end{columns} 
       
      \begin{columns}[t] % Split up the two columns wide column again
      
      \begin{column}{\sepmargin} \end{column}
        \begin{column}{\onecolwid} % The first column
	
                \vspace*{-0.9cm}
				\begin{alertblock}{\large Contact Information}
                \vspace*{-0.5cm}
					\begin{footnotesize}
					\begin{itemize}
						\item \href{mailto:email@wu.ac.at}{email@wu.ac.at}
						\item \href{http://www.example.com/}{www.example.com} - \href{www.wu.ac.at}{www.wu.ac.at}
					\end{itemize}
					\end{footnotesize}	
					
				\end{alertblock}
		    \end{column} % End of the first column
			\begin{column}{\sepwid}\end{column} % Empty spacer column
			\begin{column}{\onecolwid} % Begin a column 
              \begin{block}{\large References}
			  \vspace*{-0.5cm}
              	\nocite{*} % Insert publications even if they are not cited in the poster
					{\footnotesize
                    	%\bibliographystyle{plainurl}
						\bibliography{bibliog.bib}}
				\end{block} 
			\end{column} % End of the second column
            
			\begin{column}{\sepmargin}\end{column} % Empty spacer column
            
\end{columns} % End of all the columns in the poster


\end{frame} % End of the enclosing frame
	
\end{document}


\end{document}